\documentclass[12pt]{article}

\usepackage[T1]{fontenc}
\usepackage[utf8]{inputenc}
\usepackage{newtxtext,newtxmath}
\usepackage[margin=1in]{geometry}
\usepackage{setspace}
\usepackage{microtype}
\usepackage{xcolor}
\usepackage{enumitem}
\usepackage{hyperref}
\usepackage{fancyhdr}
\usepackage{titlesec}

\hypersetup{
  colorlinks=true,
  linkcolor=black,
  urlcolor=black,
  pdftitle={STAT 662 Presentation Speech},
  pdfauthor={Amandeep Singh}
}

\setstretch{1.12}
\setlength{\parindent}{0pt}
\setlength{\parskip}{0.55em}
\setlist[itemize]{leftmargin=1.4em, itemsep=0.2em, topsep=0.2em}

\definecolor{SpeechBlue}{HTML}{163A5F}
\definecolor{SpeechGreen}{HTML}{1E5A4B}
\definecolor{SpeechGray}{HTML}{4D5964}

\titleformat{\section}{\large\bfseries\color{SpeechBlue}}{\thesection.}{0.5em}{}
\titleformat{\subsection}{\normalsize\bfseries\color{SpeechGreen}}{}{}{}

\pagestyle{fancy}
\fancyhf{}
\fancyhead[L]{STAT 662 Presentation Speech}
\fancyhead[R]{TCGA Pan-Cancer Study}
\fancyfoot[C]{\thepage}
\setlength{\headheight}{15pt}

\newcommand{\tcga}{TCGA PanCanAtlas}

\begin{document}

\begin{center}
{\LARGE\bfseries Integrated Multivariate Analysis of \tcga\ Transcriptomic Data\par}
\vspace{0.25cm}
{\large Presenter Script for STAT 662\par}
\vspace{0.2cm}
{\normalsize Amandeep Singh\par}
\vspace{0.2cm}
{\small May 11, 2026\par}
\end{center}

\vspace{0.35cm}

{\small\color{SpeechGray}
This script is written as a print-ready speaking document for the matching Beamer presentation deck. The pacing is designed for an approximately 20-minute presentation with brief pauses for figures and transitions.
\par}

\section{Slide 1: Title}
Good morning everyone. Today I will present my STAT 662 final project, which is an integrated multivariate analysis of \tcga\ transcriptomic data.

The main motivation for this study is that modern cancer datasets are both high-dimensional and strongly structured. In particular, RNA-seq expression data contain thousands of variables per sample, but they also contain coherent biological signal that can be studied from several statistical perspectives.

Rather than focusing on only one task, such as classification or clustering alone, I built a complete workflow that studies the same pan-cancer expression dataset through four connected analyses: principal component analysis, multivariate regression, supervised cancer-type prediction, and unsupervised clustering.

After harmonization and filtering, the final cohort contains 10,055 representative tumor samples across 32 cancer types. The headline results are that the best supervised classifier achieved a held-out macro-F1 of about 0.893, while the best unsupervised clustering method achieved an adjusted Rand index of about 0.541.

So the central claim I want to support today is that pan-cancer RNA expression contains stable, clinically meaningful, and highly predictive structure, but that this structure is easier to use in supervised learning than in unsupervised discovery.

\section{Slide 2: Research Problem}
This project starts from a difficult statistical setting.

First, each tumor sample begins with more than 20,000 gene-expression measurements. That means the data are high-dimensional and strongly correlated. In this kind of setting, naive modeling often becomes unstable, and so dimension reduction and regularization are not optional, they are necessary.

Second, the cohort spans many cancer types with unequal sample sizes. That means we face both class imbalance and biological heterogeneity.

Third, some cancers are genuinely close to one another in expression space. So when a classifier confuses two cancers, or when clustering merges them, that often reflects biological similarity rather than just model failure.

Because of this, the real statistical goal is broader than simple prediction. The project asks four linked questions.

First, how much of the transcriptomic variation can be summarized by a low-dimensional latent structure?

Second, can clinical variables explain the dominant latent expression axes?

Third, how accurately can cancer type be predicted from a shared molecular representation?

And fourth, if we ignore labels completely, do tumors still group into meaningful clusters?

That is why this project is organized as one integrated pipeline instead of unrelated analyses.

\section{Slide 3: Data and Cohort Construction}
The raw data come from three main sources in the TCGA PanCanAtlas ecosystem.

The first is the pan-cancer RNA-seq expression matrix. The second is the patient-level clinical follow-up table. The third is the TCGA clinical data resource workbook.

The first major task was not modeling. It was cohort construction.

I parsed TCGA barcodes to recover patient identity and sample type. I then retained tumor-derived samples only and selected one representative tumor sample per patient, prioritizing primary solid tumors whenever multiple samples existed.

This step matters because if multiple samples from the same patient are included without adjustment, they can create pseudo-replication and distort both supervised and unsupervised results.

Starting from 11,069 manifest rows and 10,274 unique patients, the pipeline identified 10,250 representative tumor rows. After harmonization with the clinical covariates and matrix cleaning, the final cohort contained 10,055 patients across 32 cancer types.

The class distribution is broad but uneven. BRCA is the largest cohort, followed by cancers such as KIRC, UCEC, HNSC, LUAD, LGG, THCA, and LUSC.

That unevenness is one reason why later in the supervised benchmark I rely on macro-F1 rather than accuracy alone.

\section{Slide 4: Analysis Workflow}
This slide summarizes the full pipeline.

The first stage is loading and harmonizing the raw TCGA files.

The second stage is representative-sample selection and cleaning of the expression matrix.

The third stage constructs shared high-variance feature spaces, especially top-variable gene subsets.

After that, the workflow branches into four linked modules: PCA and covariance regularization, multivariate regression, supervised classification, and unsupervised clustering.

The key design choice is that these analyses are connected. They all operate on a common molecular signal rather than on unrelated feature-engineering pipelines.

That makes the final comparisons much more meaningful. For example, when supervised classification outperforms clustering, that difference reflects the presence or absence of label information rather than two completely different preprocessing choices.

For evaluation, the supervised classification module uses a stratified 80/20 train-test split with 5-fold cross-validation on the training set. PCA stability is assessed through repeated 80 percent subsampling, and clustering stability is assessed with repeated-subsample adjusted Rand index.

So throughout the project, I tried to combine performance metrics with stability diagnostics.

\section{Slide 5: PCA and Covariance Results}
The first major analysis is principal component analysis.

PCA addresses the question of whether the global transcriptomic variation can be summarized in a lower-dimensional space.

The answer is yes, but not in a trivial way.

PC1 explains 6.61 percent of the variance. The first 10 principal components explain 34.35 percent of the variance, and the first 50 explain 57.94 percent.

This means the transcriptome has substantial low-dimensional organization, but it is not a simple low-rank object with one sharp elbow. There is strong dominant structure, but also meaningful remaining variation beyond the leading components.

The second important result is stability.

I refit PCA on 10 random 80 percent subsamples and compared those solutions back to the full-data components. For PCs 1 through 5, the mean absolute loading correlations are all above 0.997, and the mean subspace similarity is 0.9995.

Those are extremely high values, and they indicate that the leading latent directions are not artifacts of one sample split. They are highly reproducible empirical summaries of the cohort.

I also examined covariance regularization because high-dimensional genomic covariance matrices are usually unstable if treated naively.

The empirical covariance matrix was strongly ill-conditioned, with a condition number of about 127,800. This justifies shrinkage and sparse inverse-covariance modeling.

Graphical Lasso produced a sparse conditional-dependence structure with biologically coherent recurring gene-gene edges, which suggests that the covariance structure is not random noise but reflects meaningful transcriptomic organization.

\section{Slide 6: Multivariate Regression}
After obtaining the latent transcriptomic structure, the next question is whether clinical variables explain that structure.

To study that, I treated the first 10 principal component scores as a multivariate response and regressed them on age, sex, pathologic stage, and cancer acronym.

Because the design matrix includes many dummy variables, I used multivariate ridge regression with cross-validated penalty selection.

This module was fit on the covariate-complete subset of 10,005 patients.

The main result is a variance-weighted cross-validated R-squared of 0.6944, with an in-sample R-squared of 0.6982. The small difference between those values suggests that the regression relation generalizes reasonably well and is not driven by major overfitting.

The most important finding, however, is the blockwise analysis.

When the cancer acronym block is removed from the design, the cross-validated R-squared drops by about 0.556. That is a very large change.

By contrast, pathologic stage contributes only a very small additional effect, and age and sex contribute very little once cancer type is already known.

So the interpretation is clear: the dominant driver of latent transcriptomic organization is cancer identity itself. Stage contributes a smaller secondary effect, while age and sex have minimal additional explanatory power after conditioning on tumor type.

This is an important bridge between PCA and classification. PCA tells us that there is strong latent structure, and regression tells us that much of that structure corresponds to cross-cancer biological differences.

\section{Slide 7: Supervised Experimental Setup}
The next part of the project is supervised cancer-type prediction.

For this task, I started from the shared top-2,000 gene expression matrix. On the training data only, I standardized the features and projected them into a 50-dimensional principal component space.

So the classifiers are not working directly on the raw 20,000-dimensional expression matrix. They are working on a lower-dimensional, cleaner, and more stable representation.

I then compared six benchmark models: Logistic Regression, Linear SVM, RBF SVM, Random Forest, Histogram Gradient Boosting, and a Multilayer Perceptron.

The data were split into 8,044 training samples and 2,011 test samples using a stratified 80/20 split, so the cancer-type proportions were preserved across training and testing.

Hyperparameters were tuned by 5-fold cross-validation on the training set only.

The key model-selection metric was macro-F1 rather than plain accuracy.

That choice matters because in a 32-class cancer prediction problem, overall accuracy can be dominated by the largest classes. Macro-F1 gives equal weight to each cancer type and is therefore a better measure of balanced classification quality.

So the question here is not simply which model has the highest accuracy, but which model performs most consistently across the full set of cancers.

\section{Slide 8: Supervised Classification Results}
The best overall model was Logistic Regression.

On the held-out test set, it achieved an accuracy of 0.9115, a macro-F1 of 0.8928, and a weighted F1 of 0.9141.

Its best cross-validated macro-F1 on the training set was 0.8727, which is close to the held-out performance and therefore supports the stability of the selection process.

One especially interesting result is that the Multilayer Perceptron had slightly higher raw accuracy, around 0.9135, but lower macro-F1, around 0.8809.

That means the neural network gained slightly on common classes but performed less evenly across the entire class set.

In this setting, I consider Logistic Regression the better model because it gives the strongest balanced performance.

There is also a broader modeling lesson here.

The strongest method is linear. Once the expression data are properly cleaned, standardized, and embedded in principal component space, a well-regularized linear decision surface is strong enough to capture most of the useful class structure.

That is a valuable practical result because linear models are easier to interpret, easier to train, and easier to deploy than more complex nonlinear alternatives.

So this part of the project shows that the pan-cancer transcriptome is highly predictive of tumor identity when supervision is available.

\section{Slide 9: Classification Diagnostics}
Aggregate metrics are useful, but they do not show where the model succeeds and where it struggles.

At the class level, the best classifier is very strong for many cancers. Twenty of the 32 cancer types achieved F1 scores of at least 0.90.

Several cancers were classified almost perfectly, including UVM, THCA, TGCT, and PCPG. Very high scores were also obtained for PRAD, BRCA, OV, LIHC, and THYM.

The harder classes were READ, UCS, ESCA, COAD, and CHOL.

Some of that difficulty is due to low support, but some of it reflects genuine biological overlap.

The confusion patterns make that point clearly.

The largest confusion pair was COAD predicted as READ, and the reverse error of READ predicted as COAD was also prominent. There were also errors between STAD and ESCA.

These are not arbitrary mistakes. They are biologically local errors between related gastrointestinal cancers.

That is actually a reassuring failure mode. If the model makes an error, we prefer it to confuse related diseases rather than distant and implausible categories.

So this slide shows two things at once: first, that the classifier performs strongly overall, and second, that the remaining errors are biologically interpretable.

\section{Slide 10: Unsupervised Clustering Results}
The next question is what happens when the labels are removed completely.

For clustering, I again used the shared top-2,000 gene space, standardized it, and reduced it to 20 principal components. I then compared three clustering methods: K-means, agglomerative clustering, and Gaussian mixture modeling.

To make the external comparison interpretable, I fixed the number of clusters at 32, matching the observed number of cancer types in the final cohort.

K-means was the strongest overall method.

It achieved a silhouette score of 0.2630, an adjusted Rand index of 0.5406, a normalized mutual information of 0.6762, and a subsample stability score of 0.8020.

Agglomerative clustering had the highest normalized mutual information, at 0.7149, but K-means performed better overall by ARI and stability.

Gaussian mixtures were somewhat weaker in this representation.

The main interpretation is that unsupervised learning does recover substantial cancer structure, but it is clearly weaker than supervised classification.

That is expected. Supervised models are given label information and are allowed to optimize directly for prediction, whereas clustering is not.

Still, an ARI above 0.5 and strong stability indicate that known tumor identities are major organizing principles of the expression manifold.

So the clustering analysis supports the supervised results, but from a weaker information setting.

\section{Slide 11: Integrated Interpretation}
At this point, the separate analyses can be combined into one statistical story.

PCA tells us that the transcriptome has strong low-dimensional structure.

The PCA stability analysis tells us that those latent directions are highly reproducible.

Regression tells us that the dominant driver of that latent structure is cancer identity, while stage contributes a smaller secondary effect.

Supervised classification shows that this shared molecular representation is highly predictive when labels are available.

Unsupervised clustering shows that even without labels, the data still recover substantial biological grouping.

The most important comparison is between supervised and unsupervised performance.

The best classifier reaches a macro-F1 of about 0.893, whereas the best clustering method reaches an ARI of about 0.541.

That gap tells us that tumor identity is strongly encoded in the expression data, but it is not perfectly separable without supervision.

So the overall message is not just that the dataset contains signal. It is that the same molecular representation supports explanation, inference, prediction, and discovery in a coherent way.

If I had to summarize the project in one sentence, I would say that pan-cancer RNA expression contains stable and clinically useful latent structure, but supervision makes much more efficient use of that structure than unsupervised learning does.

\section{Slide 12: Limitations and Next Steps}
Like any study, this one has limitations.

First, it is an expression-only analysis. It does not integrate mutation, methylation, copy-number, or survival endpoints.

Second, it uses one representative tumor sample per patient, which is appropriate statistically, but does not capture within-patient heterogeneity.

Third, clustering was evaluated at a fixed 32-cluster setting, which is useful for comparison but may not match the natural number of molecular subgroups.

Fourth, the regression module uses a focused clinical covariate panel rather than a broader clinical design.

These limitations suggest natural extensions.

One direction is multimodal genomic integration, combining expression with other molecular data types.

Another is shifting from broad cancer-type prediction toward subtype discovery within individual cancers.

A third direction is extending the regression component toward survival modeling or other clinically important outcomes.

And finally, the learned representation could be validated on an external cohort to evaluate generalizability beyond TCGA.

So I see this project as a strong statistical baseline and a reusable framework, rather than a final endpoint.

\section{Slide 13: Closing}
To conclude, this project implements a complete multivariate analysis pipeline for \tcga\ pan-cancer transcriptomic data.

Starting from raw RNA-seq and clinical files, it constructs a patient-level tumor cohort, builds shared high-variance feature spaces, and uses them consistently across PCA, covariance modeling, regression, supervised classification, clustering, and report generation.

Numerically, the main findings are the following.

First, the transcriptome has strong and highly stable low-dimensional structure.

Second, cancer type is the dominant source of explainable latent variation.

Third, Logistic Regression provides the best balanced supervised classification performance, with a held-out macro-F1 of 0.8928.

And fourth, K-means provides meaningful but weaker unsupervised recovery of tumor identity, with an adjusted Rand index of 0.5406.

So the central statistical conclusion of this project is that pan-cancer transcriptomic data are rich enough to support a coherent multivariate analysis program in which dimension reduction, covariance estimation, regression, classification, and clustering each illuminate a different aspect of the same biological structure.

Thank you. I am happy to take questions.

\section{Short Q\&A Notes}
If a question comes up after the presentation, these are useful short responses to keep in mind.

\begin{itemize}
  \item Why did you use macro-F1 instead of accuracy? Because the class sizes are uneven, and macro-F1 gives balanced weight to all 32 cancers.
  \item Why did Logistic Regression outperform more complex models? Because the expression space becomes highly structured after standardization and PCA, so a regularized linear boundary is already very effective.
  \item Why is clustering weaker than classification? Because clustering has no access to labels, while classification is directly optimized for separating the observed cancer types.
  \item Why did the regression use fewer than 10,055 patients? Because the regression model requires complete covariate information, so only the covariate-complete subset of 10,005 patients enters that analysis.
  \item What is the main biological takeaway? Cancer identity is the dominant driver of latent transcriptomic variation, and related cancers remain the hardest to separate.
\end{itemize}

\end{document}
